\documentclass[12pt]{article}
\usepackage[a4paper,margin=1in]{geometry}
\usepackage[utf8]{inputenc}
\usepackage[T1]{fontenc}
\usepackage{amsmath,amssymb}
\usepackage{textcomp}
\usepackage{parskip}
\begin{document}
\begin{center}
{\Large \textbf{2000 STEP II}}
\end{center}

\subsection*{Page 1}
\begin{flushleft}
STEP II, 2000 2\par
\bigskip\textbf{Section A: Pure Mathematics}\par
\bigskip\textbf{Question 1.} A number of the form 1/N , where N is an integer greater than 1, is called a unit fraction.\par
Noting that\par
1 1 1 1 1 1\par
\[= + and = + ,\]
2 3 6 3 4 12\par
guess a general result of the form\par
1 1 1\par
\[= + ( )\]
N a b ∗\par
and hence prove that any unit fraction can be expressed as the sum of two distinct unit\par
fractions.\par
By writing ( ) in the form\par
∗ 2\par
\[(a N )(b N ) = N\]
- -\par
and by considering the factors of N 2, show that if N is prime, then there is only one way of\par
expressing 1/N as the sum of two distinct unit fractions.\par
Prove similarly that any fraction of the form 2/N , where N is prime number greater than 2,\par
can be expressed uniquely as the sum of two distinct unit fractions.\par
\bigskip\textbf{Question 2.} Prove that if (x a)2 is a factor of the polynomial p (x), then p 0 (a) = 0. Prove a corresponding\par
result if (x\par
a)4-\par
is a factor of p (x).\par
-\par
Given that the polynomial\par
6 5 4 3 2\par
x + 4x 5x 40x 40x + 32x + k\par
- - -\par
has a factor of the form (x a) , find k.\par
-\par
\bigskip\textbf{Question 3.} The lengths of the sides BC, CA, AB of the triangle ABC are denoted by a, b, c, respectively.\par
Given that\par
\[b = 8 + ²1, c = 3 + ²2, A = \pi/3 + ²3,\]
where ²1, ²2, and ²3 are small, show that a 7 + η, where η = (13 ²1 2 ²2 + 24\sqrt3 ²3)/14.\par
\approx -\par
Given now that\par
\[²1 6 2 10 -3 , ²2 6 4 9 10 -2 , ²3 6 \sqrt\{3\} 10 -3 ,\]
\[| | \times | | · \times | | \times\]
find the range of possible values of η.\par
\end{flushleft}
\subsection*{Page 2}
\begin{flushleft}
STEP II, 2000 3\par
\bigskip\textbf{Question 4.} Prove that\par
\[(cos θ + i sin θ)(cos φ + i sin φ) = cos(θ + φ) + i sin(θ + φ)\]
and that, for every positive integer n,\par
n\par
\[(cos θ + i sin θ) = cos nθ + i sin nθ.\]
By considering (5 i)2(1 + i), or otherwise, prove that\par
-\par
\[arctan(7/17) + 2 arctan(1/5) = \pi/4.\]
Prove also that\par
\[3 arctan(1/4) + arctan(1/20) + arctan(1/1985) = \pi/4.\]
[Note that arctan θ is another notation for tan\par
-1\par
θ.]\par
5 It is required to approximate a given function f (x), over the interval 0 6 x 6 1, by the linear\par
function λx, where λ is chosen to minimise\par
1 2\par
f (x) λx dx.\par
\[Z^0 ³ - ´\]
Show that\par
\[λ = 3 xf (x) dx.\]
\[Z^0\]
The residual error, R, of this approximation process is such that\par
1 2\par
\[R = f (x) λx dx.\]
\[Z^0 ³ - ´\]
Show that\par
1 2 1\par
2 2\par
\[R = f (x) dx λ .\]
\[Z^0 ³ ´ - 3\]
Given now that f (x) = sin(\pix/n), show that (i) for large n, λ \pi/n and (ii) lim n\to\infty R = 0.\par
\approx\par
Explain why, prior to any calculation, these results are to be expected.\par
\[[You may assume that, when θ is small, sin θ θ θ3/6 and cos θ 1 θ2/2.]\]
\approx - \approx -\par
\end{flushleft}
\subsection*{Page 3}
\begin{flushleft}
STEP II, 2000 4\par
\bigskip\textbf{Question 6.} Show that\par
\[2t 1 t^2 1 + cos θ\]
\[sin θ = , cos θ = - , = tan(\pi/2 θ/2),\]
\[1 + t^2 1 + t^2 sin θ -\]
\[where t = tan(θ/2).\]
\[Use the substitution t = tan(θ/2) to show that, for 0 < α < \pi/2,\]
\[\pi\]
2 1 α\par
\[dθ = ,\]
\[Z^0 1 + cos α sin θ sin α\]
and deduce a similar result for\par
\[\pi\]
2 1\par
dθ.\par
\[Z^0 1 + sin α cos θ\]
\bigskip\textbf{Question 7.} The line l has vector equation r = λs, where\par
\[s = (cos θ + \sqrt\{3\} ) i + (\sqrt\{2\} sin θ) j + (cos θ \sqrt\{3\} ) k\]
-\par
and λ is a scalar parameter. Find an expression for the angle between l and the line\par
\[r = µ(a i + b j + c k). Show that there is a line m through the origin such that, whatever\]
\[the value of θ, the acute angle between l and m is \pi/6.\]
A plane has equation x z = 4\sqrt3. The line l meets this plane at P . Show that, as θ varies,\par
-\par
P describes a circle, with its centre on m. Find the radius of this circle.\par
\bigskip\textbf{Question 8.} (i) Let y be the solution of the differential equation\par
dy\par
+ 4x e\par
\[-x^2\]
(y + 3)\par
\[2 = 0 (x > 0),\]
dx\par
\[that satisfies the condition y = 6 when x = 0. Find y in terms of x and show that y 1\]
\to\par
as x .\par
\to \infty\par
\textbf{(ii)} Let y be any solution of the differential equation\par
\[dy x e 6x2 (y + 3) 1-k = 0 (x > 0).\]
dx -\par
Find a value of k such that, as x , e\par
\[-3x2\]
y tends to a finite non-zero limit, which\par
\to \infty\par
you should determine.\par
\end{flushleft}
\subsection*{Page 4}
\begin{flushleft}
STEP II, 2000 5\par
\bigskip\textbf{Section B: Mechanics}\par
\bigskip\textbf{Question 9.} In an aerobatics display, Jane and Karen jump from a great height and go through a pe-\par
riod of free fall before opening their parachutes. While in free fall at speed v, Jane experi-\par
ences air resistance kv per unit mass but Karen, who spread-eagles, experiences air resistance\par
\[kv + (2k2/g)v^2 per unit mass. Show that Jane's speed can never reach g/k. Obtain the\]
corresponding result for Karen.\par
Jane opens her parachute when her speed is g/(3k). Show that she has then been in free fall\par
for time k\par
-1\par
\[ln(3/2).\]
Karen also opens her parachute when her speed is g/(3k). Find the time she has then been\par
in free fall.\par
\bigskip\textbf{Question 10.} A long light inextensible string passes over a fixed smooth light pulley. A particle of mass 4 kg\par
is attached to one end A of this string and the other end is attached to a second smooth light\par
pulley. A long light inextensible string BC passes over the second pulley and has a particle\par
of mass 2 kg attached at B and a particle of mass of 1 kg attached at C. The system is\par
held in equilibrium in a vertical plane. The string BC is then released from rest. Find the\par
accelerations of the two moving particles.\par
After T seconds, the end A is released so that all three particles are now moving in a vertical\par
plane. Find the accelerations of A, B and C in this second phase of the motion. Find also, in\par
terms of g and T , the speed of A when B has moved through a total distance of 0.6gT 2 metres.\par
\bigskip\textbf{Question 11.} The string AP has a natural length of 1 5metres and modulus of elasticity equal to 5g newtons.\par
·\par
The end A is attached to the ceiling of a room of height 2 5 metres and a particle of mass\par
·\par
0 5 kg is attached to the end P . The end P is released from rest at a point 0 5 metres above\par
· ·\par
the floor and vertically below A. Show that the string becomes slack, but that P does not\par
reach the ceiling.\par
Show also that while the string is in tension, P executes simple harmonic motion, and that\par
the time in seconds that elapses from the instant when P is released to the instant when P\par
first returns to its original position is\par
1 1\par
8 2 3 2\par
\[+ \pi arccos(3/7) .\]
µ 3g ¶ µ 5g ¶ -\par
³ ´\par
[Note that arccos x is another notation for cos\par
-1\par
x.]\par
\end{flushleft}
\subsection*{Page 5}
\begin{flushleft}
STEP II, 2000 6\par
\bigskip\textbf{Section C: Probability and Statistics}\par
12 Tabulated values of Φ( ), the cumulative distribution function of a standard normal variable,\par
·\par
should not be used in this question.\par
Henry the commuter lives in Cambridge and his working day starts at his office in London at\par
0900. He catches the 0715 train to King's Cross with probability p, or the 0720 to Liverpool\par
Street with probability 1 p. Measured in minutes, journey times for the first train are\par
-\par
N (55, 25) and for the second are N (65, 16). Journey times from King's Cross and Liverpool\par
Street to his office are N (30, 144) and N (25, 9), respectively. Show that Henry is more likely\par
to be late for work if he catches the first train.\par
\[Henry makes M journeys, where M is large. Writing A for 1 Φ(20/13) and B for 1 Φ(2),\]
- -\par
find, in terms of A, B, M and p, the expected number, L, of times that Henry will be late\par
and show that for all possible values of p,\par
BM 6 L 6 AM.\par
\[Henry noted that in 3/5 of the occasions when he was late, he had caught the King's Cross\]
train. Obtain an estimate of p in terms of A and B.\par
[A random variable is said to be N µ, σ2 if it has a normal distribution with mean µ and\par
variance σ2.] ¡ ¢\par
\end{flushleft}
\subsection*{Page 6}
\begin{flushleft}
STEP II, 2000 7\par
\bigskip\textbf{Question 13.} A group of biologists attempts to estimate the magnitude, N , of an island population of voles\par
(Microtus agrestis). Accordingly, the biologists capture a random sample of 200 voles, mark\par
them and release them. A second random sample of 200 voles is then taken of which 11 are\par
found to be marked. Show that the probability, p , of this occurrence is given by\par
N\par
(N 200)!\par
-\par
\[p = k ³ ´ ,\]
N\par
N !(N 389)!\par
-\par
where k is independent of N .\par
The biologists then estimate N by calculating the value of N for which p is a maximum.\par
N\par
Find this estimate.\par
All unmarked voles in the second sample are marked and then the entire sample is released.\par
Subsequently a third random sample of 200 voles is taken. Write down the probability that this\par
sample contains exactly j marked voles, leaving your answer in terms of binomial coefficients.\par
Deduce that\par
200\par
389 3247 3636\par
\[= .\]
µ j ¶µ200 j¶ µ 200 ¶\par
\[Xj=0 -\]
\bigskip\textbf{Question 14.} The random variables X1, X2, . . . , X2n+1 are independently and uniformly distributed on\par
the interval 0 6 x 6 1. The random variable Y is defined to be the median of X1, X2, . . . ,\par
\[X2n+1. Given that the probability density function of Y is g (y), where\]
kyn(1 y) n if 0 6 y 6 1\par
\[g (y) = -\]
½ 0 otherwise,\par
use the result\par
r!s!\par
\[yr(1 y) s dy =\]
\[Z^0 - (r + s + 1)!\]
to show that k = (2n + 1)!/(n!) , and evaluate E(Y ) and Var (Y ). Hence show that, for any\par
given positive number d, the inequality\par
\[P Y 1/2 < d/\sqrt\{n\} < P X¯ 1/2 < d/\sqrt\{n\}\]
| - | | - |\par
¡ ¢ ¡ ¢\par
\[holds provided n is large enough, where X¯ is the mean of X^1, X^2, . . . , X2n+1.\]
[You may assume that Y and X¯ are normally distributed for large n.]\par
\end{flushleft}
\end{document}
